\documentclass[11pt, a4paper]{article}

% ====== Required Packages ======
\usepackage[utf8]{inputenc}
\usepackage[T1]{fontenc}
\usepackage{lmodern}
\usepackage{microtype}
\usepackage{amsmath, amssymb, amsfonts}
\usepackage{geometry}
\usepackage{booktabs}
\usepackage{tabularx}
\usepackage{array}
\usepackage[table]{xcolor}
\usepackage{hyperref}
\usepackage{enumitem}
\usepackage{fancyhdr}
\usepackage{listings}
\usepackage{titlesec}

% ====== Page Geometry ======
\geometry{top=2.5cm, bottom=2.5cm, left=2.5cm, right=2.5cm, headheight=14pt}

% ====== Colors ======
\definecolor{accentblue}{RGB}{30, 80, 160}
\definecolor{passgreen}{RGB}{198, 239, 206}
\definecolor{warnyellow}{RGB}{255, 235, 156}
\definecolor{warnred}{RGB}{255, 199, 206}
\definecolor{tableheader}{RGB}{230, 236, 245}
\definecolor{codebg}{RGB}{246, 248, 250}
\definecolor{codekw}{RGB}{170, 13, 145}
\definecolor{codecmt}{RGB}{0, 110, 60}
\definecolor{codestr}{RGB}{196, 26, 22}

% ====== Section formatting ======
\titleformat{\section}{\large\bfseries\color{accentblue}}{\thesection}{1em}{}
\titleformat{\subsection}{\normalsize\bfseries}{\thesubsection}{1em}{}

% ====== Hyperlink Setup ======
\hypersetup{
  colorlinks=true,
  linkcolor=accentblue,
  urlcolor=accentblue,
  pdftitle={DAMPS-MMHCL Wave 1 Compliance Audit and M1.5 Sweep Driver},
}

% ====== Listings (Python) ======
\lstdefinestyle{py}{
  language=Python,
  backgroundcolor=\color{codebg},
  basicstyle=\ttfamily\scriptsize,
  keywordstyle=\color{codekw}\bfseries,
  commentstyle=\color{codecmt}\itshape,
  stringstyle=\color{codestr},
  numbers=left,
  numberstyle=\tiny\color{gray},
  numbersep=8pt,
  showstringspaces=false,
  breaklines=true,
  frame=single,
  rulecolor=\color{gray!40},
  tabsize=4,
  columns=fullflexible,
  literate={->}{{$\rightarrow$}}1 {>=}{{$\geq$}}1 {<=}{{$\leq$}}1
}

% ====== Custom marks ======
\newcommand{\cmark}{\textcolor{green!60!black}{\textbf{\checkmark}}}
\newcommand{\warnmark}{\textcolor{orange!80!black}{\textbf{!}}}

\newcolumntype{Y}{>{\centering\arraybackslash}X}
\newcolumntype{Z}{>{\raggedright\arraybackslash}X}

% ====== Header / Footer ======
\pagestyle{fancy}
\fancyhf{}
\fancyhead[L]{\small\textit{DAMPS-MMHCL Wave 1 --- LogQ Compliance Audit}}
\fancyhead[R]{\small\thepage}
\fancyfoot[C]{\small Phase 2 Execution Verification Report}
\renewcommand{\headrulewidth}{0.4pt}
\renewcommand{\footrulewidth}{0.4pt}

\title{%
  \textbf{Wave 1 Implementation Audit: Decoupled LogQ Deployment (1B $\rightarrow$ 1A)}\\[0.4em]
  \large DAMPS-MMHCL Phase 2 --- Section 8.2 Compliance Verification\\[0.2em]
  \normalsize Code-Level vs.\ Execution-Level Completion Assessment
}
\author{DAMPS-MMHCL Research Team}
\date{\today}

% ============================================================
\begin{document}
% ============================================================

\maketitle
\tableofcontents
\newpage

% ------------------------------------------------------------
\section{Executive Summary}
% ------------------------------------------------------------

The 1B/1A source code achieves \textbf{100\% conformance} to the design
specification --- including the subtle safety details (fail-fast behaviour,
item-branch-only application, and float32 overflow protection). However, while the
implementation is complete at the code level, \textbf{Wave 1 is not yet 100\%
``complete'' at the execution level}, because the M1.5 milestone has never actually
been run. The following sections detail the gap.

\begin{center}
\renewcommand{\arraystretch}{1.4}
\begin{tabular}{lcc}
\toprule
\textbf{Completion Dimension} & \textbf{Level} & \textbf{Status} \\
\midrule
\cellcolor{passgreen}Code-level (1B + 1A implementation) & $\sim$100\% &
  \cmark~Spec-conformant, incl.\ safety details \\
\cellcolor{warnred}Execution-level (M1.5 milestone) & Not met &
  \warnmark~No run, no logs, no sweep \\
\bottomrule
\end{tabular}
\end{center}

% ------------------------------------------------------------
\section{Issues Preventing Wave 1 From Reaching 100\% ``Completion''}
% ------------------------------------------------------------

\subsection{Issue 1 --- The training notebook never enables LogQ (the largest gap)}

The notebook \texttt{Local\_Random\_seeds\_train\_mmhcl\_clothing} currently runs
only the rev44 Phase~1 configuration (\texttt{apc\_off\_combined}). A full inspection
of every \texttt{train.py} invocation reveals only nine \texttt{-{}-damps\_*} flags
(\texttt{apc}, \texttt{avrf}, \texttt{imcf}, \texttt{fft}, \texttt{routing},
\texttt{momentum}, \texttt{prior}, \texttt{categories}, \texttt{warmup}) and
\textbf{not a single \texttt{-{}-enable\_logq} line} (zero occurrences of the token
\texttt{logq} anywhere in the notebook). This means the Wave~1 code has been merged
but \textbf{never executed} --- the M1.5 milestone has no data.

\subsection{Issue 2 --- No evidence of a 10/10 green test run}

The M1.5 README mandates that ``all tests must be green before merging 1A,'' but the
repository contains no CI artifact or pytest log. The output of
\texttt{pytest tests/test\_popularity\_prior.py -v} must be attached as evidence.

\subsection{Issue 3 --- No backward-compatibility smoke test}

The README requires running one epoch with \texttt{-{}-enable\_logq 0} and then
performing a bit-for-bit \texttt{diff} against rev45. No log demonstrating that this
check was performed has been found.

\subsection{Issue 4 --- The \texttt{logq\_scale} sweep has not been run}

The specification stresses that \texttt{logq\_scale} $\in \{0.05, 0.1, 0.3, 1.0\}$
must be swept \emph{before} trusting the literal specification, because
\texttt{scale}\,$=1.0$ produces logits of magnitude $\sim 40$, which overflow.
This sweep is a precondition of the M1.5 gate and is currently absent.

\subsection{Issue 5 --- A minor technical risk inside the LogQ code path}

In the \texttt{use\_logq} branch, \texttt{col\_slice} is passed as the \emph{entire}
\texttt{indices} for every chunk, whereas \texttt{refl\_sim}/\texttt{between\_sim} are
computed over \texttt{z1[mask]} (the current chunk only). It must be confirmed that
\texttt{log\_q\_term[col\_slice]} is applied correctly along the column axis (the full
$N$-item dimension) and is not misaligned when \texttt{batch\_size}\,$<$\,\texttt{num\_nodes}.
The current logic appears correct (columns $=$ the full item set), but a small numerical
test should verify the loss value when $N >$ batch size.

% ------------------------------------------------------------
\section{Compliance Matrix (Code Level)}
% ------------------------------------------------------------

\renewcommand{\arraystretch}{1.35}
\begin{table}[h]
\centering
\caption{Code-level conformance of the 1B + 1A implementation against rev54 §8.2.}
\begin{tabularx}{\textwidth}{|Z|Z|c|}
\hline
\rowcolor{tableheader}
\textbf{Spec requirement (rev54 §8.2)} & \textbf{Actual implementation} & \textbf{Status} \\
\hline
Split 1B (estimator) from 1A (loss) into two PRs &
  \texttt{damps/popularity\_prior.py} is pure; loss patch is separate &
  \cellcolor{passgreen}\cmark \\
\hline
$q(i)$ Laplace: $(n_i+\beta)/(\sum_j n_j + |I|\cdot\beta)$ &
  \texttt{compute\_log\_q} laplace mode matches the formula &
  \cellcolor{passgreen}\cmark \\
\hline
Support raw vs.\ Laplace (warning §3.1) &
  three modes available: laplace / raw / sqrt &
  \cellcolor{passgreen}\cmark \\
\hline
Cache $+$ invalidation when the train split changes &
  SHA1 signature over (mode, $\beta$, $n_{items}$, $\sum$counts) &
  \cellcolor{passgreen}\cmark \\
\hline
Independent unit tests for the estimator &
  10 tests (T1--T10) $+$ 1 slow test &
  \cellcolor{passgreen}\cmark \\
\hline
\texttt{enable\_logq=False} $\Rightarrow$ bit-for-bit identical to rev45 &
  \texttt{f\_simple} branch preserved; default flag $=0$ &
  \cellcolor{passgreen}\cmark \\
\hline
Fail-fast when the log\_q buffer is all zeros &
  \texttt{RuntimeError} if \texttt{log\_q.abs().sum()==0} &
  \cellcolor{passgreen}\cmark \\
\hline
LogQ applies to the item branch only, NOT the user branch &
  item \texttt{apply\_logq=True}, user \texttt{apply\_logq=False} &
  \cellcolor{passgreen}\cmark \\
\hline
Two extra safety hyperparameters (\texttt{logq\_scale}, \texttt{logq\_clip}) against float32 overflow &
  present $+$ clip $\pm 5.0$, scale sweepable &
  \cellcolor{passgreen}\cmark \\
\hline
Five CLI flags (\texttt{enable\_logq}, \texttt{logq\_mode/beta/scale/clip}) &
  added to \texttt{parser.py} &
  \cellcolor{passgreen}\cmark \\
\hline
\end{tabularx}
\end{table}

% ------------------------------------------------------------
\section{Summary}
% ------------------------------------------------------------

\begin{itemize}[leftmargin=*]
  \item \textbf{Code level: $\sim$100\%} --- conformant to the spec, including the
    safety details.
  \item \textbf{Wave 1 execution level: not met} --- because M1.5 has not been run:
    the notebook does not enable \texttt{-{}-enable\_logq}, there is no pytest log, no
    backward-compatibility smoke test, and no scale sweep.
\end{itemize}

\noindent\textbf{Remaining work to reach 100\% Wave~1:}
\begin{enumerate}[leftmargin=*]
  \item Add a notebook cell/variant that passes
    \texttt{-{}-enable\_logq 1 -{}-logq\_mode laplace -{}-logq\_beta 1.0 -{}-logq\_scale \{sweep\}}.
  \item Attach the 10/10 pytest log.
  \item Run the backward-compatibility diff.
  \item Run the 24-run sweep and report the M1.5 gate
    (R@20 $\geq 0.0925$, with no seed below $0.0890$).
\end{enumerate}

% ------------------------------------------------------------
\section{Appendix: M1.5 Sweep Driver Notebook Cell}
% ------------------------------------------------------------

The cell below executes the M1.5 sweep over the matrix
\texttt{logq\_scale} $\times$ \texttt{logq\_mode} $\times$ \texttt{seed}. It calls
\texttt{train.py} once per configuration, parses the \texttt{BEST\_Test\_*} metrics
from each run, aggregates per-(mode, scale) means and standard deviations across
seeds, and evaluates the M1.5 gate. Paste it as a new cell at the end of the
notebook (it reuses \texttt{DAMPS\_DIR}, \texttt{dataset}, and \texttt{PYTHON\_EXE}
defined in the earlier cells).

\begin{lstlisting}[style=py, caption={M1.5 sweep driver: logq\_scale $\times$ logq\_mode $\times$ seed.}]
# =====================================================================
# Section 8.2 -- Wave 1 / M1.5 sweep driver
#   Matrix:  logq_scale x logq_mode x seed   (calls train.py per config)
#   Gate:    mean Recall@20 >= 0.0925  AND  every seed >= 0.0890
# =====================================================================
from __future__ import annotations

import os
import re
import sys
import json
import time
import random
import subprocess
from typing import Any

# ---- 0. Resolve project dirs (reuse earlier-cell globals if present) ----
try:
    _DAMPS_DIR = DAMPS_DIR
except NameError:
    _DAMPS_DIR = os.getcwd()
try:
    _PYTHON = PYTHON_EXE
except NameError:
    _PYTHON = sys.executable
try:
    _DATASET = dataset
except NameError:
    _DATASET = "clothing"

if os.path.normpath(os.getcwd()) != os.path.normpath(_DAMPS_DIR):
    os.chdir(_DAMPS_DIR)

# ---- 1. Sweep matrix definition --------------------------------------
LOGQ_SCALES: list[float] = [0.05, 0.1, 0.3, 1.0]   # spec-mandated scale sweep
LOGQ_MODES:  list[str]   = ["laplace"]             # add "raw","sqrt" to widen
N_SEEDS: int = 5                                   # 4 scales x 1 mode x 5 = 20 runs
LOGQ_BETA: float = 1.0
LOGQ_CLIP: float = 5.0

# M1.5 acceptance gate (rev54 Section 8.2 / Section 6)
GATE_MEAN_RECALL: float = 0.0925   # mean R@20 must reach this
GATE_MIN_RECALL:  float = 0.0890   # NO seed may fall below this (rollback gate)

# Reuse the seed list from the Phase 1 sweep if available, else generate.
try:
    SEEDS: list[int] = list(seeds)[:N_SEEDS]          # noqa: F821
    if len(SEEDS) < N_SEEDS:
        raise NameError
except NameError:
    random.seed(2026)
    SEEDS = [random.randint(1, 10_000) for _ in range(N_SEEDS)]
print(f"[M1.5] seeds = {SEEDS}")

# ---- 2. Fixed rev45 apc_off_combined backbone flags ------------------
#   Identical to the Phase-1 'combined' variant so the ONLY delta is LogQ.
BACKBONE_FLAGS: list[str] = [
    "--dataset",                 _DATASET,
    "--temperature",             "0.3",
    "--damps_apc",               "0",
    "--damps_avrf",              "0",
    "--damps_imcf",              "1",
    "--damps_permutation_fft",   "1",
    "--damps_soft_routing",      "1",
    "--damps_momentum",          "1",
    "--damps_data_driven_prior", "1",
    "--damps_warmup_epochs",     "0",
]

# ---- 3. Metric parser (matches BEST_Test_* in train.py text logs) ----
_PAT = {
    "recall@20": re.compile(r"BEST_Test_Recall@20[^0-9]*([0-9]*\.?[0-9]+)"),
    "ndcg@20":   re.compile(r"BEST_Test_NDCG@20[^0-9]*([0-9]*\.?[0-9]+)"),
}

def _parse_metrics(stdout: str) -> dict[str, float]:
    out: dict[str, float] = {}
    for key, pat in _PAT.items():
        hits = pat.findall(stdout)
        if hits:
            out[key] = float(hits[-1])   # last = best-checkpoint snapshot
    return out

# ---- 4. Run one (mode, scale, seed) configuration --------------------
def run_one(mode: str, scale: float, seed: int) -> dict[str, Any]:
    cmd: list[str] = [
        _PYTHON, "train.py",
        *BACKBONE_FLAGS,
        "--seed",        str(seed),
        # ---- Wave 1 LogQ activation (the delta vs rev45) ----
        "--enable_logq", "1",
        "--logq_mode",   mode,
        "--logq_beta",   str(LOGQ_BETA),
        "--logq_scale",  str(scale),
        "--logq_clip",   str(LOGQ_CLIP),
    ]
    t0 = time.time()
    proc = subprocess.run(cmd, capture_output=True, text=True, cwd=_DAMPS_DIR)
    dt = time.time() - t0
    log = proc.stdout + "\n" + proc.stderr
    metrics = _parse_metrics(log)
    ok = ("recall@20" in metrics) and (proc.returncode == 0)
    if not ok:
        print(f"  [WARN] mode={mode} scale={scale} seed={seed} "
              f"rc={proc.returncode} -> last 400 chars:\n{log[-400:]}")
    return {
        "mode": mode, "scale": scale, "seed": seed,
        "recall@20": metrics.get("recall@20"),
        "ndcg@20":   metrics.get("ndcg@20"),
        "runtime_s": round(dt, 1), "ok": ok,
    }

# ---- 5. Execute the full matrix --------------------------------------
all_runs: list[dict[str, Any]] = []
total = len(LOGQ_MODES) * len(LOGQ_SCALES) * len(SEEDS)
i = 0
print("=" * 72)
print(f"M1.5 SWEEP: {total} runs "
      f"({len(LOGQ_MODES)} mode x {len(LOGQ_SCALES)} scale x {len(SEEDS)} seed)")
print("=" * 72)
for mode in LOGQ_MODES:
    for scale in LOGQ_SCALES:
        for seed in SEEDS:
            i += 1
            print(f"[{i}/{total}] mode={mode} scale={scale} seed={seed} ...",
                  flush=True)
            all_runs.append(run_one(mode, scale, seed))

# ---- 6. Aggregate per (mode, scale) and evaluate the gate ------------
def _mean(xs: list[float]) -> float:
    return sum(xs) / len(xs) if xs else float("nan")

def _std(xs: list[float]) -> float:
    if len(xs) < 2:
        return 0.0
    m = _mean(xs)
    return (sum((x - m) ** 2 for x in xs) / (len(xs) - 1)) ** 0.5

summary: list[dict[str, Any]] = []
for mode in LOGQ_MODES:
    for scale in LOGQ_SCALES:
        rs = [r["recall@20"] for r in all_runs
              if r["mode"] == mode and r["scale"] == scale
              and r["recall@20"] is not None]
        ns = [r["ndcg@20"] for r in all_runs
              if r["mode"] == mode and r["scale"] == scale
              and r["ndcg@20"] is not None]
        if not rs:
            continue
        mean_r, min_r = _mean(rs), min(rs)
        passed = (mean_r >= GATE_MEAN_RECALL) and (min_r >= GATE_MIN_RECALL)
        summary.append({
            "mode": mode, "scale": scale, "n": len(rs),
            "recall_mean": round(mean_r, 4), "recall_std": round(_std(rs), 4),
            "recall_min": round(min_r, 4),
            "ndcg_mean": round(_mean(ns), 4), "ndcg_std": round(_std(ns), 4),
            "M1.5_pass": passed,
        })

# ---- 7. Report -------------------------------------------------------
print("\n" + "=" * 72)
print("M1.5 GATE REPORT  (mean R@20 >= 0.0925  AND  min R@20 >= 0.0890)")
print("=" * 72)
hdr = f"{'mode':<9}{'scale':>7}{'n':>4}{'R@20 mean':>12}{'+/-':>9}" \
      f"{'R@20 min':>11}{'NDCG mean':>11}{'PASS':>7}"
print(hdr)
print("-" * len(hdr))
for s in summary:
    print(f"{s['mode']:<9}{s['scale']:>7}{s['n']:>4}"
          f"{s['recall_mean']:>12.4f}{s['recall_std']:>9.4f}"
          f"{s['recall_min']:>11.4f}{s['ndcg_mean']:>11.4f}"
          f"{('YES' if s['M1.5_pass'] else 'no'):>7}")

any_pass = any(s["M1.5_pass"] for s in summary)
print("\n[M1.5 VERDICT]",
      "PASS -- at least one (mode,scale) cleared the gate; "
      "promote that config to Wave 2 (M1 / SimGCL)."
      if any_pass else
      "FAIL -- no (mode,scale) cleared the gate; "
      "inspect logq_scale/clip before proceeding.")

# ---- 8. Persist raw + summary for the audit trail --------------------
out_dir = os.path.abspath(os.path.join("..", _DATASET))
os.makedirs(out_dir, exist_ok=True)
with open(os.path.join(out_dir, "m15_logq_sweep_runs.json"), "w") as f:
    json.dump(all_runs, f, indent=2)
with open(os.path.join(out_dir, "m15_logq_sweep_summary.json"), "w") as f:
    json.dump(summary, f, indent=2)
print(f"\nSaved -> {out_dir}/m15_logq_sweep_runs.json")
print(f"Saved -> {out_dir}/m15_logq_sweep_summary.json")
\end{lstlisting}

\noindent\textbf{Usage notes.}
\begin{itemize}[leftmargin=*]
  \item The cell assumes the earlier environment-setup cell has defined
    \texttt{DAMPS\_DIR}, \texttt{PYTHON\_EXE}, and \texttt{dataset}; if not, it falls
    back to sensible defaults.
  \item The \texttt{BACKBONE\_FLAGS} list reproduces the rev45 \texttt{apc\_off\_combined}
    backbone so that the \emph{only} delta versus the validated baseline is the LogQ
    activation, preserving causal attributability for the M1.5 ablation.
  \item To widen the sweep into the full 24-run matrix mentioned in the audit, add
    \texttt{"raw"} (and optionally \texttt{"sqrt"}) to \texttt{LOGQ\_MODES}; the gate
    logic and reporting scale automatically.
  \item Adjust the regular expressions in \texttt{\_PAT} if your \texttt{train.py}
    prints metric keys under a different label than \texttt{BEST\_Test\_Recall@20}.
\end{itemize}

\end{document}
